\begin{example*}
The real plane $\mathbb{R}^2$ carries the usual Euclidean metric
\[
  d(p,q)
  =
  \sqrt{(p_1-q_1)^2 + (p_2-q_2)^2}.
\]
In Lean, we represent a point of the plane as an ordered pair
\[
  \texttt{p : Real × Real}.
\]
The first coordinate is \texttt{p.1}, and the second coordinate is
\texttt{p.2}. Thus the Euclidean distance formula is written as
\[
  \texttt{Real.sqrt ((p.1 - q.1) \char`^ 2 + (p.2 - q.2) \char`^ 2)}.
\]

The Lean definition begins:
\[
\begin{gathered}
  \texttt{noncomputable def realPlaneEuclideanMetric : Metric (Real × Real) where}\\
  \texttt{\ \ distance p q :=}\\
  \texttt{\ \ \ \ Real.sqrt ((p.1 - q.1) \char`^ 2 + (p.2 - q.2) \char`^ 2).}
\end{gathered}
\]
It is marked \texttt{noncomputable} because \texttt{Real.sqrt} is a
mathematical real-valued function, not executable computation data.

The first metric axiom is nonnegativity. Mathematically, for arbitrary
$p,q \in \mathbb{R}^2$, we must show
\[
  0 \le
  \sqrt{(p_1-q_1)^2 + (p_2-q_2)^2}.
\]
This follows immediately from the basic fact that real square roots are
nonnegative.

In Lean, Mathlib provides this theorem as \texttt{Real.sqrt\_nonneg}. The proof
is:
\[
\begin{gathered}
  \texttt{distance\_nonnegative := by}\\
  \texttt{\ \ intro p q}\\
  \texttt{\ \ exact Real.sqrt\_nonneg \_.}
\end{gathered}
\]
The command \texttt{intro p q} introduces arbitrary points
\texttt{p q : Real × Real}. After this, the goal is the nonnegativity of the
specific square root expression defining the distance from \texttt{p} to
\texttt{q}. The theorem \texttt{Real.sqrt\_nonneg} says that
\[
  0 \le \sqrt{r}
\]
for any real number $r$. The underscore in
\[
  \texttt{Real.sqrt\_nonneg \_}
\]
asks Lean to infer the real number under the square root from the current goal.
Here Lean infers
\[
  r = (p.1-q.1)^2 + (p.2-q.2)^2.
\]
So this proof is not using any special property of sums of squares yet; it only
uses the general nonnegativity of real square roots.

The next axiom is identity of indiscernibles:
\[
  d(p,q) = 0 \iff p = q.
\]
For the Euclidean distance, this becomes
\[
  \sqrt{(p_1-q_1)^2 + (p_2-q_2)^2} = 0 \iff p = q.
\]
As usual, \texttt{constructor} splits the proof into two directions. The
forward direction assumes that the distance is zero and proves that the points
are equal. The reverse direction assumes that the points are equal and proves
that the distance is zero.

The forward direction starts from
\[
  \texttt{hypothesis : Real.sqrt ((p.1 - q.1) \char`^ 2 + (p.2 - q.2) \char`^ 2) = 0}.
\]
First we prove that the expression under the square root is nonnegative:
\[
\begin{gathered}
  \texttt{have h\_nonneg :}\\
  \texttt{\ \ 0 <= (p.1 - q.1) \char`^ 2 + (p.2 - q.2) \char`^ 2 := by}\\
  \texttt{\ \ positivity.}
\end{gathered}
\]
The tactic \texttt{positivity} proves goals that follow from standard
nonnegativity facts, such as squares being nonnegative and sums of
nonnegative terms being nonnegative.

Then \texttt{Real.sqrt\_eq\_zero h\_nonneg} rewrites the statement that the
square root is zero into the statement that the nonnegative expression under
the square root is zero. After this rewrite, Lean can use \texttt{nlinarith} to
show that each square must be zero:
\[
  (p.1-q.1)^2 = 0,
  \qquad
  (p.2-q.2)^2 = 0.
\]
The theorem \texttt{sq\_eq\_zero\_iff} turns each squared equation into a
linear equation, and \texttt{sub\_eq\_zero} turns
\[
  p_i-q_i = 0
\]
into
\[
  p_i=q_i.
\]
Finally, \texttt{ext} proves equality of ordered pairs by proving equality of
their two coordinates.

The reverse direction is simpler. If \texttt{hypothesis : p = q}, then rewriting
by \texttt{hypothesis} changes the distance from \texttt{p} to \texttt{q} into
the distance from \texttt{q} to itself. The command \texttt{ring\_nf} reduces
the expression under the square root to \texttt{0}, and \texttt{Real.sqrt\_zero}
finishes the proof.

The Lean proof is:

\begin{verbatim}
distance_eq_zero_iff_equal := by
  intro p q
  constructor
  -- =>
  . intro hypothesis
    have sum_of_squares_eq_zero :
        (p.1 - q.1) ^ 2 + (p.2 - q.2) ^ 2 = 0 := by
      have h_nonneg : 0 <= (p.1 - q.1) ^ 2 + (p.2 - q.2) ^ 2 := by
        positivity
      rwa [Real.sqrt_eq_zero h_nonneg] at hypothesis
    ext
    . have h_sq : (p.1 - q.1) ^ 2 = 0 := by
        nlinarith
      rwa [sq_eq_zero_iff, sub_eq_zero] at h_sq
    . have h_sq : (p.2 - q.2) ^ 2 = 0 := by
        nlinarith
      rwa [sq_eq_zero_iff, sub_eq_zero] at h_sq
  -- <=
  . intro hypothesis
    rw [hypothesis]
    ring_nf
    exact Real.sqrt_zero
\end{verbatim}

The symmetry axiom asks us to prove
\[
  d(x,y) = d(y,x).
\]
For the Euclidean metric on \(\mathbb{R}^2\), this means proving
\[
  \sqrt{(x_1-y_1)^2 + (x_2-y_2)^2}
  =
  \sqrt{(y_1-x_1)^2 + (y_2-x_2)^2}.
\]
The proof is:
\[
\begin{gathered}
  \texttt{distance\_symmetric := by}\\
  \texttt{\ \ intro x y}\\
  \texttt{\ \ congr 1}\\
  \texttt{\ \ ring.}
\end{gathered}
\]
The command \texttt{intro x y} introduces arbitrary points
\texttt{x y : Real × Real}. After this, the goal is the equality of the two
square-root expressions above.

The command \texttt{congr 1} uses congruence. In this proof, both sides have
the same outer function, namely \texttt{Real.sqrt}. The goal has the form
\[
  \texttt{Real.sqrt A = Real.sqrt B}.
\]
To prove this, it is enough to prove
\[
  \texttt{A = B}.
\]
So \texttt{congr 1} strips away one matching outer layer of the expression and
leaves the new goal
\[
  (x.1-y.1)^2 + (x.2-y.2)^2
  =
  (y.1-x.1)^2 + (y.2-x.2)^2.
\]
This is now a pure algebraic identity. It follows from the fact that
\[
  (a-b)^2 = (b-a)^2.
\]
The tactic \texttt{ring} proves this kind of identity by normalizing both sides
as ring expressions. After normalization, both sides have the same canonical
form, so Lean recognizes them as equal.

Thus the proof has two conceptual steps: first reduce equality of square roots
to equality of the expressions under the square roots, and then prove that
inner equality by algebra.

For the triangle inequality, we can prove the result by connecting our
coordinate definition on \(\mathbb{R}^2\) to Mathlib's built-in Euclidean
space. This is an interoperation proof: our metric is still our own
\texttt{Metric (Real × Real)}, but we temporarily translate points into
Mathlib's \texttt{EuclideanSpace} so that we can use Mathlib's already-proved
triangle inequality for its distance function.

The extra import is:
\[
  \texttt{import Mathlib.Analysis.InnerProductSpace.PiL2}.
\]
This file provides Mathlib's finite-dimensional Euclidean spaces, their
distance function, and facts such as \texttt{dist\_triangle}. In particular, we
use \texttt{EuclideanSpace ℝ (Fin 2)}, the Euclidean space of real-valued
functions on the two-element type \texttt{Fin 2}.

Our points are represented as pairs \texttt{ℝ × ℝ}, while Mathlib's Euclidean
space represents a point as a length-two coordinate vector. The conversion is:
\[
\begin{gathered}
  \texttt{noncomputable def toE (p : ℝ × ℝ) : EuclideanSpace ℝ (Fin 2) :=}\\
  \texttt{\ \ (EuclideanSpace.equiv (Fin 2) ℝ).symm ![p.1, p.2].}
\end{gathered}
\]
The vector notation \texttt{![p.1, p.2]} builds the length-two coordinate
vector whose first coordinate is \texttt{p.1} and whose second coordinate is
\texttt{p.2}. The equivalence \texttt{EuclideanSpace.equiv (Fin 2) ℝ} relates
Mathlib's Euclidean space to these coordinate vectors. We use \texttt{.symm}
because we are moving from coordinates into Euclidean space.

The bridge lemma says that our square-root formula is exactly Mathlib's
Euclidean distance after this conversion:
\[
\begin{gathered}
  \texttt{lemma distance\_eq\_euclidean (p q : ℝ × ℝ) :}\\
  \texttt{\ \ Real.sqrt ((p.1 - q.1) \char`^ 2 + (p.2 - q.2) \char`^ 2)}\\
  \texttt{\ \ \ \ = dist (toE p) (toE q) := by}\\
  \texttt{\ \ rw [EuclideanSpace.dist\_eq, Fin.sum\_univ\_two]}\\
  \texttt{\ \ simp [toE, EuclideanSpace.equiv, Real.dist\_eq, sq\_abs].}
\end{gathered}
\]
The proof unfolds Mathlib's Euclidean distance formula with
\texttt{EuclideanSpace.dist\_eq}, rewrites the sum over \texttt{Fin 2} as the
two coordinate terms using \texttt{Fin.sum\_univ\_two}, and then simplifies the
coordinate conversion. The theorem \texttt{Real.dist\_eq} rewrites real
distance as absolute value, and \texttt{sq\_abs} removes absolute values under
squares.

With this bridge lemma, the triangle inequality proof becomes:
\[
\begin{gathered}
  \texttt{distance\_triangle := by}\\
  \texttt{\ \ intro x y z}\\
  \texttt{\ \ rw [distance\_eq\_euclidean, distance\_eq\_euclidean,}\\
  \texttt{\ \ \ \ distance\_eq\_euclidean]}\\
  \texttt{\ \ exact dist\_triangle (toE x) (toE y) (toE z).}
\end{gathered}
\]
After \texttt{intro x y z}, the goal is the triangle inequality for our
coordinate square-root formula. The three rewrites replace each occurrence of
our formula with Mathlib's distance between the converted points:
\[
  d(x,z),\qquad d(x,y),\qquad d(y,z).
\]
The goal is then exactly Mathlib's triangle inequality:
\[
  \texttt{dist (toE x) (toE z)}
  \le
  \texttt{dist (toE x) (toE y) + dist (toE y) (toE z)}.
\]
The theorem \texttt{dist\_triangle} proves this for any pseudo-metric/metric
space known to Mathlib, including \texttt{EuclideanSpace ℝ (Fin 2)}.

This is a useful shortcut, but it is also a little bit of a cheat for the
exercise: instead of proving the Euclidean triangle inequality directly from
algebra and square roots, we proved that our formula matches Mathlib's
Euclidean distance and then reused Mathlib's triangle inequality.
\end{example*}
